\begin{proof}
\begin{enumerate}[label=\arabic*.,leftmargin=*]
\item \emph{The holomorphic family.}
Let $Q$ be a disc with $\overline Q\subset\mathcal P$.
On an ordinary lifted base chart use the period matrix
\eqref{wb:pshift:eq:periods}, jointly holomorphic in $(z,c)$.
Its real determinant is $\Im\tau\,D_c$, and its inverse is
uniformly bounded on compact base-parameter sets. Thus the lattice
action is free and properly discontinuous there. A fixed closed
real lattice parallelotope gives compact surjective representatives
over each compact subset of the base chart times $Q$.
The quotient is holomorphic and its projection to $Q$ is a submersion.
The covariance convention of \cref{I:def:affine-covariance} is
\[
 R_g(z)\Pi(z,c)=\Pi(gz,c)A_g,
 \qquad M_g=A_g^{-1},\qquad
 \Pi(gz,c)=R_g(z)\Pi(z,c)M_g.
\]
Consequently the base deck quotient is jointly holomorphic; its
proper discontinuity follows from that of the lifted base action.

At a finite root chart the family action is
\begin{equation}
 (z,c,\zeta)\longmapsto
 \left(g_jz,c,R_{g_j}(z)\zeta+
             \Pi(g_jz,c)\frac{v_j}{m_j}\right),
 \qquad (m_1,m_2)=(3,4).
\label{pub:I:eq:PD7}
\end{equation}
In the real marking it is $(z,c,u)\mapsto(g_jz,c,A_ju+v_j/m_j)$.
The invariant-covector freeness test in \cref{fam:src:A.1.2}
is independent of $c$, so these quotients are local biholomorphisms
and submersions to $Q$. The same parallelotope proves properness
over the finite base disc times $Q$. Their punctured connectors are
\begin{equation}
 \frac{\log s_j(z)}{2\pi i}\Pi(z,c)v_j.
\label{pub:I:eq:PD8}
\end{equation}
A logarithm change adds a period, so this is a jointly holomorphic
torus section. Writing $\sigma_j=(\log s_j)/(2\pi i)\,\Pi(z,c)v_j$, the
relations $A_jv_j=v_j$ and
$\log s_j(g_jz)=\log s_j(z)-2\pi i/m_j$ give
\[
 R_{g_j}(z)\sigma_j(z,c)-\sigma_j(g_jz,c)
       =\Pi(g_jz,c)v_j/m_j.
\]
Thus translation by $\sigma_j$ intertwines the finite affine action
with the middle linear action, with the specified root coordinate.

\item \emph{Uniform cusp charts and properness.}
On the fixed unimodular fan write the full cusp period matrix as
\begin{equation}
 C(t_c,c)=\begin{pmatrix}6\mu(t_c)&h(t_c)\\
 b^{\rm part}(t_c)-h(t_c)+c&\mu(t_c)\end{pmatrix},
 \qquad B_0=\begin{pmatrix}0&1\\-1&0\end{pmatrix}.
\label{pub:I:eq:PD9}
\end{equation}
Choose a closed cusp disc $|t_c|\le r_1$ on which these functions
are holomorphic, and set
\begin{equation}
 K_Q=\max_{|t_c|\le r_1,\ c\in\overline Q}
               \|-2\pi\Im C(t_c,c)\|,
 \qquad 0<\epsilon<\min(r_1,1),\quad
               |\log\epsilon|>2K_Q.
\label{pub:I:eq:PD10}
\end{equation}
The lattice acts on the toric family times $Q$ by
\begin{equation}
 (x,t_c,c)\longmapsto
 \bigl(e^{2\pi iC(t_c,c)\lambda}t_c^{B_0\lambda}x,t_c,c\bigr),
 \qquad\lambda\in\mathbb Z^2.
\label{pub:I:eq:PD11}
\end{equation}
The fan translation and nowhere-zero torus multiplier extend this
map across every toric stratum. With Euclidean operator norm,
$B_0^{\mathsf t}B_0=I$, and
\begin{equation}
 B_{t_c,c}=B_0+\frac{-2\pi\Im C(t_c,c)}{\log|t_c|},
 \quad \|B_{t_c,c}-B_0\|<\tfrac12,
 \quad \tfrac12\|\lambda\|<\|B_{t_c,c}\lambda\|
                    \le\tfrac32\|\lambda\|\quad(\lambda\ne0).
\label{pub:I:eq:PD12}
\end{equation}
These inequalities exclude stabilizers on the dense torus.
For a height-one triangle, use its three dual unimodular cone
coordinates $z_0,z_1,z_2$, with $t_c=z_0z_1z_2$.
In the open chart $|z_i|<2$, $|t_c|<\epsilon$, the normalized
logarithmic barycentric coefficients have sum one and are bounded
below by $-\log2/|\log\epsilon|$. Thus the logarithmic positions
in each fixed chart lie in a bounded set, uniformly in $c$.
For a fixed pair of charts, \eqref{pub:I:eq:PD12} bounds the finitely many
$\lambda$ which can translate one into the other. At the central
fibre a nonzero $B_0\lambda$ moves every bounded fan cell.
These charts cover the entire toric space, including its boundary.
On the dense torus the duality identity is
\[
 \log|z_i|=\log|t_c|\,\theta_i^\sigma(y),
\]
where $\theta_i^\sigma$ are the barycentric coordinates of the
height-one triangle. A triangle containing $y$ therefore gives
$|z_i|\le1$. For a boundary point in the orbit of the cone over a
cell $P$, the coordinates belonging to vertices of $P$ vanish.
The remaining characters are dual to the image cones of triangles
containing $P$ in the quotient lattice by $\operatorname{span}(P,1)$.
Those cones cover the quotient real space: lift a vector, add a
sufficiently small positive multiple of it to a relative interior
point of $(P,1)$, and rescale to height one. Local finiteness places
that point in a triangle containing $P$. Applied to the negative
logarithm of the orbit coordinates, this gives the same closed
unit-polydisc bound. Thus the open radius-two charts cover the
boundary as well. Every compact set has a finite subcover.
The chart-pair estimate applies to an intersection meeting the
boundary because that intersection is open and meets the dense
torus. It therefore proves freeness and proper discontinuity on
the entire toric family.

For properness, fix $\epsilon_1<\epsilon$ and let $K_0$ be the
finite union of closed unit polydiscs for height-one triangles
meeting the closed radius-$\sqrt2$ ball, intersected with
$|t_c|\le\epsilon_1$. For $t_c\ne0$, choose $\lambda$ with
$\|B_{t_c,c}^{-1}y+\lambda\|_\infty\le1/2$.
The translated logarithmic position has norm at most
$3/(2\sqrt2)<\sqrt2$. A triangle containing it has nonnegative
barycentric coordinates of sum one, so each corresponding cone
coordinate has modulus at most one. This gives a representative
in $K_0$. At $t_c=0$, translate the height-one component and use
the six closed toric bidiscs covering $E_0\simeq\mathrm dP_6$.
The multiplier $\exp(2\pi iC(0,c)B_0v)$ is included in this
translation; the bidiscs cover the resulting whole compact component.
For each compact $K\subset Q$, the compact set $K_0\times K$
therefore maps onto the full inverse image of
$\{|t_c|\le\epsilon_1\}\times K$. This proves properness.
The quotient charts are open subsets of a smooth toric variety
times $Q$, so projection to $Q$ is a submersion.

\item Exponentiating the logarithmic coordinates gives the cusp
connector with period identity $Z=sB_0+C(t_c,c)$.
Glue it and \eqref{pub:I:eq:PD8} to the middle family. The distinguished base
discs are disjoint. Points with distinct $(t,c)$ values separate
on the base; points with the same value lie in one Hausdorff local
piece. The resulting total space is therefore Hausdorff and smooth.
The local proper maps give a proper map $F$ to $B\times Q$:
cover a compact subset by finitely many relatively compact sets
inside the displayed base pieces. Compactness of $B$ proves
properness of $\pi=\operatorname{pr}_QF$.
The fibres and overlaps are the specified $X_c$ and $f_c$.
These local families agree on smaller cusp radii and glue over
$\mathcal P$.
\item \emph{Local maps.}
Let $n\in2\mathbb Z$. Define the following rational columns,
integral cocycles, and holomorphic translation vectors:
\begin{equation}
\begin{aligned}
 r_1&=-\frac n3(\widehat u+\widehat w),&
 \ell_1&=-n\widehat w,&
 b_1(z)&=\Pi_c(z)r_1=-\frac n3(\tau(z)+1,\mu(z))^{\mathsf t},\\
 r_2&= \frac n4(\widehat u+\widehat w),&
 \ell_2&=\frac n2\widehat u,&
 b_2(z)&=\Pi_c(z)r_2=\frac n4(\tau(z)+1,\mu(z))^{\mathsf t}.
\end{aligned}
\hypertarget{tag-03CH}{}\label{wb:pshift:eq:local-corrections}
\end{equation}
Direct multiplication of the complete matrices in \eqref{wb:pshift:eq:finite-data} gives
\begin{equation}
 (I-A_j)r_j=\frac{n\epsilon_j}{m_j}\widehat\delta+\ell_j.
 \hypertarget{tag-03CI}{}\label{wb:pshift:eq:local-matrix-identity}
\end{equation}
For $j=1$ the four entries are $(0,0,-n,n/3)$, and for $j=2$
they are $(0,n/2,0,-n/4)$. Since both $r_j$ and
$n\epsilon_j\widehat\delta/m_j+\ell_j$ lie in $\ker\gamma$,
the period covariance gives
\[
 b_j(g_jz)-R_{g_j}(z)b_j(z)
 =\Pi_{c+n}(g_jz)
       \left(\frac{n\epsilon_j}{m_j}\widehat\delta+\ell_j\right).
\]
This is the deck-equivariance identity for translation by $b_j$
between the finite fillings at $c$ and $c+n$. Its overlap with
the fixed middle identification is as follows.

Keep the marked linearizing coordinate, radius, and phase
$s_j(g_jz)=e^{-2\pi i/m_j}s_j(z)$. On a logarithm branch put
\begin{equation}
 w_j(z)=\frac{\log s_j(z)}{2\pi i},\qquad
 a_j(z)=b_j(z)+n\epsilon_j w_j(z)e.
 \hypertarget{tag-03CA}{}\label{wb:pshift:eq:finite-middle-sections}
\end{equation}
Changing the logarithm adds $n\epsilon_j\widehat\delta$ under the
period map, so $a_j$ is branch independent as a torus section.
For the specified negative generator, $w_j(g_jz)=w_j(z)-1/m_j$.
Equations \eqref{wb:pshift:eq:covariance} and \eqref{wb:pshift:eq:local-matrix-identity} consequently give the exact equality
\begin{equation}
 a_j(g_jz)=R_{g_j}(z)a_j(z)+\Pi_c(g_jz)\ell_j.
 \hypertarget{tag-03CG}{}\label{wb:pshift:eq:local-affine-cocycle}
\end{equation}
In fact the marked finite overlap maps are translations by
$\sigma_{j,c}=w_j\Pi_c v_j$, and
\begin{equation}
                  \sigma_{j,c+n}-\sigma_{j,c}
                     =n\epsilon_jw_je.
 \hypertarget{tag-03BX}{}\label{wb:pshift:eq:actual-overlap-difference}
\end{equation}
Thus $a_j$ is exactly the middle translation induced by the local
translation $b_j$, with the marked gluing maps retained.

At the cusp retain $s=\tau-h$, $e^{2\pi is}=t_c$, $s(g_0z)=s(z)-1$,
and the complete period matrix
\begin{equation}
 Z_c=sB_0+C_c(t_c),\quad
 B_0=\begin{pmatrix}0&1\\-1&0\end{pmatrix},\quad
 C_c=\begin{pmatrix}6\mu&h\\ b^{\mathrm{part}}-h+c&\mu\end{pmatrix}.
 \hypertarget{tag-03C5}{}\label{wb:pshift:eq:cusp-periods}
\end{equation}
For $\lambda=(\lambda_1,\lambda_2)^{\mathsf t}\in\mathbb Z^2$,
$C_{c+n}\lambda-C_c\lambda=(0,n\lambda_1)^{\mathsf t}$. Therefore
\begin{equation}
 \exp(2\pi iC_{c+n}\lambda)t_c^{B_0\lambda}
       =\exp(2\pi iC_c\lambda)t_c^{B_0\lambda}.
 \hypertarget{tag-03C3}{}\label{wb:pshift:eq:cusp-actions-identical}
\end{equation}
The entire cusp actions, including their nonconstant holomorphic units,
agree under the integer-shift identification.

The affine cocycle convention in \eqref{wb:pshift:eq:local-affine-cocycle} has multiplication law
\begin{equation}
                    \ell_{gh}=\ell_g+A_g\ell_h.
 \hypertarget{tag-03C0}{}\label{wb:pshift:eq:cocycle-law}
\end{equation}
Since $g_0=(g_1g_2)^{-1}$, it gives
\begin{equation}
\begin{aligned}
 \ell_1+A_1\ell_2
   &=-\frac{3n}{2}\widehat w+\frac n2\widehat\delta,\\
 \ell_0=-A_{g_0}(\ell_1+A_1\ell_2)
   &=\frac{3n}{2}\widehat w-\frac n2\widehat\delta.
\end{aligned}
\hypertarget{tag-03C4}{}\label{wb:pshift:eq:cusp-cocycle}
\end{equation}
Here $A_{g_0}$ fixes both toric lattice vectors, as is also immediate
from $A_{g_0}=\left(\begin{smallmatrix}I&0\\B_0&I\end{smallmatrix}\right)$.
Put
\begin{equation}
 q_n=(-3n/2,n/2)^{\mathsf t}\in\mathbb Z^2,
                  \qquad a_0(z)=s(z)q_n.
 \hypertarget{tag-03C6}{}\label{wb:pshift:eq:cusp-translation}
\end{equation}
Then $a_0(g_0z)-a_0(z)=-q_n=\Pi_c(g_0z)\ell_0$, exactly as
required by \eqref{wb:pshift:eq:local-affine-cocycle}. Exponentiating gives the map
$(x,t_c)\mapsto(t_c^{q_n}x,t_c)$ on the dense cusp torus.
It extends holomorphically and invertibly to the marked infinite fan:
the integral map
\[
 (y,y_3)\longmapsto(y+y_3q_n,y_3)
\]
translates each height-one triangle by the integer vector $q_n$ and
therefore permutes the cones, with inverse given by $-q_n$.
It commutes with the cusp quotient action in \eqref{wb:pshift:eq:cusp-actions-identical}. This proves extension across every cusp stratum.
\item \emph{The global additive correction.}
Let $\widetilde B^\circ$ be
the universal cover of the three-punctured base. The fundamental group
of $B^\circ$ is free on the marked two negative peripheral generators. Pull back
the marked $\Pi_c,R$ to this cover. Assign $\ell_1,\ell_2$ from
\eqref{wb:pshift:eq:local-corrections} to its free generators and extend by \eqref{wb:pshift:eq:cocycle-law}. Freeness makes
this an integral cocycle; its third peripheral value is
precisely \eqref{wb:pshift:eq:cusp-cocycle}. The maps
\begin{equation}
 (z,y)\longmapsto
       (gz,R_g(z)y+\Pi_c(gz)\ell_g)
 \hypertarget{tag-03BY}{}\label{wb:pshift:eq:affine-bundle-action}
\end{equation}
form an affine action. To check the action law, covariance gives
$R_g(hz)\Pi_c(hz)=\Pi_c(ghz)A_g$; substituting this equality
in the composition gives exactly \eqref{wb:pshift:eq:cocycle-law}.
The quotient is an affine holomorphic bundle over $B^\circ$, since
the deck action on the base is free and properly discontinuous and
one may use evenly covered base opens as affine charts.

On a finite punctured disc, \eqref{wb:pshift:eq:local-affine-cocycle} says that $a_j$ is a section of
this affine bundle. On the cusp punctured disc, \eqref{wb:pshift:eq:cusp-cocycle}--\eqref{wb:pshift:eq:cusp-translation} say
the same for $a_0$. Define a sheaf $\mathscr H_n$ on $B$ as follows.
Away from the special points it is the sheaf of sections of \eqref{wb:pshift:eq:affine-bundle-action}.
Near $p_j$, a section is required to have the form $a_j+\eta_j$,
where $\eta_j:\Delta_j\to\mathbb C^2$ is holomorphic and obeys
$\eta_j(g_jz)=R_{g_j}(z)\eta_j(z)$. Near the cusp it is required
to have the form $a_0+\eta_0$, with $\eta_0(t_c)$ holomorphic on
the full cusp disc. These are compatible definitions on the punctured
overlaps by the equal affine cocycles just proved. They also prove
local nonemptiness, including at all three special points.

The differences of these local sections form a vector sheaf
$\mathscr V$. Its ordinary sections have covariance $\eta(gz)=R_g\eta(z)$;
at the two finite points and the cusp its extensions are precisely
the holomorphic extensions specified above. Projection to the first
coordinate gives an exact sequence on the marked compact base
\begin{equation}
 0\longrightarrow\mathcal O_B e\longrightarrow\mathscr V
        \longrightarrow\mathcal L\longrightarrow0,
                  \qquad \mathcal L\simeq\mathcal O_B(-p_0).
 \hypertarget{tag-03CP}{}\label{wb:pshift:eq:vector-sheaf}
\end{equation}
The homogeneous first-coordinate line is determined by its transition factors. Indeed
the first diagonal entries in \eqref{wb:pshift:eq:covariance} are its factors $-1/\tau$ and
$1/\tau$, and at the cusp both extensions require ordinary
holomorphicity. The marked function
\[
 \Theta=\frac{(E_4\circ\tau)^2(E_6\circ\tau)^{1/2}}
                 {\Delta_{\mathrm{mod}}\circ\tau}
\]
has those exact factors, orders two and one in $s_1,s_2$, and a
simple pole in $t_c$. Dividing by $\Theta$ identifies the ordinary
and finite stalks with $\mathcal O_B$, and the cusp stalk with
$\mathfrak m_{p_0}$. This recovers the displayed $\mathcal L$
with its fixed cusp condition. 
For the finite stalks in \eqref{wb:pshift:eq:vector-sheaf}, use the full matrix $R_g$. Given an equivariant first component
$u(z)$, begin with the vector $(u(z),0)^{\mathsf t}$ and average it
under the finite group action on vector-valued holomorphic functions.
The resulting vector is equivariant, and its first component is
still $u$. Thus projection is surjective there. Its kernel is
exactly the invariant second-component functions, which are
holomorphic functions of $s_j^{m_j}=t_j$, multiplied by $e$.
The ordinary and cusp stalks are immediate. This also proves that
$\mathscr V$ is locally free of rank two: lift a local frame of
$\mathcal L$ and join it to $e$.

Both $H^1(B,\mathcal O_B)$ and
$H^1(B,\mathcal O_B(-p_0))$ vanish: on the two standard affine
charts of $\mathbb P^1$, split a holomorphic Laurent series on
their overlap into its nonnegative and negative powers. For
$\mathcal O_B$ these are the two coboundary contributions, and
for $\mathcal O_B(-p_0)$ its transition factor shifts the second
contribution by one, still leaving no missing exponent. The positive
part converges on the full first affine chart and the transformed
negative part on the full second affine chart. This proves the two
vanishings, equivalently the usual additive Cousin calculation on
these charts. The exact sequence \eqref{wb:pshift:eq:vector-sheaf} now gives
\begin{equation}
                             H^1(B,\mathscr V)=0.
 \hypertarget{tag-03CJ}{}\label{wb:pshift:eq:no-analytic-obstruction}
\end{equation}

Choose
local sections $h_\alpha$ of $\mathscr H_n$, taking the displayed
$a_j,a_0$ on sufficiently small special discs and affine charts
elsewhere. Their differences $d_{\alpha\beta}=h_\beta-h_\alpha$
are a Cech $1$-cocycle with values in $\mathscr V$. Project it to
$\mathcal L$ and solve that additive cocycle, using its just proved
$H^1$ vanishing. Lift the resulting local $\mathcal L$ sections to
$\mathscr V$ by the local lifting construction in \eqref{wb:pshift:eq:vector-sheaf}. Subtracting
their coboundary from $d$ leaves a scalar $\mathcal O_Be$ cocycle.
Solve that scalar additive cocycle by its $H^1$ vanishing. This
produces sections $v_\alpha$ with
$d_{\alpha\beta}=v_\beta-v_\alpha$. The local sections
$h_\alpha-v_\alpha$ agree and define a global section
$a\in\mathscr H_n(B)$. It gives holomorphic vectors $\eta_j,\eta_0$ with
\begin{equation}
 a=a_j+\eta_j\quad(j=1,2),\qquad a=a_0+\eta_0
                              \quad\hbox{at the cusp}.
 \hypertarget{tag-03CE}{}\label{wb:pshift:eq:global-correction}
\end{equation}
Any choices made in the additive solutions give maps with the same
asserted properties. Their differences as sections of this affine
bundle lie in $H^0(B,\mathscr V)=\mathbb C e$: this follows from
\eqref{wb:pshift:eq:vector-sheaf}, $H^0(B,\mathcal L)=0$, and $H^0(B,\mathcal O_B)=\mathbb C$.

Projecting $a$ modulo $\Pi_c\Lambda$ gives a holomorphic section of
the middle torus family because every $\ell_g$ in \eqref{wb:pshift:eq:affine-bundle-action} is integral.
The three local maps are therefore
\begin{equation}
\begin{aligned}
 F_{\mathrm{mid}}[z,\zeta]&=[z,\zeta+a(z)],\\
 F_j[z,\zeta]&=[z,\zeta+b_j(z)+\eta_j(z)]\quad(j=1,2),\\
 F_0[x,t_c]&=[\exp(2\pi i\eta_0(t_c))t_c^{q_n}x,t_c].
\end{aligned}
\hypertarget{tag-03CD}{}\label{wb:pshift:eq:full-isomorphism}
\end{equation}
The first descends by \eqref{wb:pshift:eq:affine-bundle-action} and equality of the integer-shift
lattices. The finite maps descend by
\eqref{wb:pshift:eq:local-matrix-identity}, the finite
deck-equivariance identity above, and the exact covariance of $\eta_j$. The cusp map descends by \eqref{wb:pshift:eq:cusp-actions-identical}; its
holomorphic-unit factor extends by torus multiplication on every
toric chart, and its integer factor extends by the fan map in
\eqref{wb:pshift:eq:cusp-translation}. Each map is a biholomorphism, with inverse given by the
negative translation or the inverse torus multiplier and fan map.

On a finite overlap, write the marked gluing as
$h_{j,c}(\zeta)=\zeta+\sigma_{j,c}$. Then \eqref{wb:pshift:eq:actual-overlap-difference}, \eqref{wb:pshift:eq:global-correction}, and
\eqref{wb:pshift:eq:full-isomorphism} give the equality of torus maps
\begin{equation}
                     h_{j,c+n}F_j=F_{\mathrm{mid}}h_{j,c}.
 \hypertarget{tag-03CC}{}\label{wb:pshift:eq:finite-overlap-identity}
\end{equation}
Both sides add
$\sigma_{j,c}+a=\sigma_{j,c+n}+b_j+\eta_j$.
On the cusp overlap the marked map is exponentiation of $\zeta$;
\eqref{wb:pshift:eq:cusp-translation} and \eqref{wb:pshift:eq:global-correction} give
$\exp(2\pi ia)=t_c^{q_n}\exp(2\pi i\eta_0)$, exactly the
third line of \eqref{wb:pshift:eq:full-isomorphism}. Thus that overlap also agrees, including
all nonconstant period units. The maps and their inverses glue to
a global biholomorphism $X_c\to X_{c+n}$ over $B$.

The choices can be made coherently for every parameter and every even
shift. Both initial cocycle columns in \eqref{wb:pshift:eq:local-corrections} have $\gamma=0$.
Because every $A_g$ preserves $\gamma$, \eqref{wb:pshift:eq:cocycle-law} shows that
$\gamma(\ell_g)=0$ for every word. Therefore the affine term
$\Pi_c(gz)\ell_g$ in \eqref{wb:pshift:eq:affine-bundle-action} is independent of $c$. The displayed
$a_j,a_0$, and the vector sheaf $\mathscr V$ are also independent
of $c$. Perform the additive construction once for $n=2$, keeping
one resulting section $a^{(2)}$. For $n=2k$ choose
$a^{(2k)}=k a^{(2)}$ and multiply its local correcting vectors by
$k$. All their affine cocycles and local formulas are then exactly
those for $n=2k$. The resulting maps satisfy the full group law
\begin{equation}
 F_{2h,c+2k}\circ F_{2k,c}=F_{2(h+k),c},\qquad
                 F_{0,c}=\mathrm{id}_{X_c}.
 \hypertarget{tag-03C1}{}\label{wb:pshift:eq:coherent-shift-action}
\end{equation}
This follows by adding the translation vectors in the first
two lines of \eqref{wb:pshift:eq:full-isomorphism}, and by multiplying their units and adding the
integer exponents in its third line. The formulas are jointly
holomorphic in $c$, since their corrections are independent of it.
\item The local families constructed above are proper and holomorphic.
The coherent $2\mathbb Z$ action is properly
discontinuous because its action on $\mathcal P$ is.
Local inverse logarithms for $w=e^{-\pi ic}$ identify its quotient
with the proper local families just proved.
Proposition~\ref{fam:src:A.2.1} gives pairwise non-biholomorphic fibres.
\end{enumerate}
\end{proof}
